\section{Introduction}\label{sec:introduction}
Soft interval constraints specify a price for violating a preference while
leaving selected physical limits immutable. In predictive control, a preferred
state range and an actuator capacity are not interchangeable constraints.
In energy scheduling, an upper grid departure and a lower comfort departure
can have different units and consequences. The numerical problem must state
which rows may yield and which objective measures that choice before a solver
comparison has a definite meaning. Convex quadratic objectives with such
interval penalties provide a useful, transparent setting for that analysis.

An ordinary slack lift expresses many of these prices exactly, but enlarges
the mathematical system. Generalized slack projections replace selected slack
variables with proximal row operations and can preserve the original principal
dimensions~\cite{gsp,lew}. That observation motivates a computational hypothesis:
an original-row implementation might reduce the cost of a complete optimization
operation. It does not prove the hypothesis. A proximal operation can itself
be expensive, sparse fill depends on more than dimension, and preparation,
validation and recovery contribute to the measured operation.

We study two related original-row updates. The first applies separate
nonnegative linear--quadratic prices to lower and upper interval departures.
Its scalar proximal minimizer follows from exterior quadratic stationarity
and endpoint subgradient conditions. The second prices directional disjoint
group peaks. A box-residual identity under coordinate contractions reduces
its positive diagonal-metric proximal operation to a weighted scalar threshold.
We state the precise fixed-metric and relaxation scope, preserve hard rows
separately, and test outputs in original physical coordinates. Existing
side-weighted active-set formulations, directional row splitting and compact
epigraphs already express the same objectives~\cite{daqp,gsp}. Consequently,
neither unequal side costs nor directional expressibility is a priority claim.

The two computational evaluations are separately registered synthetic studies.
The original study contains 456 instances, with 168 primary paired observations
from 24 distinct held-out problems and seven repetitions across six seed clusters.
Its median cold end-to-end reduction is \OriginalEstimate\%, with 95\% limits
\OriginalLower\% and \OriginalUpper\%. The grouped study contains 210 instances
and a complete 144-instance held-out population. Its primary score is
\PrimaryEstimate\%, with limits \PrimaryLower\% and \PrimaryUpper\%, including
\PrimaryFailed{} failed pairs. Both predicted 10\% improvements are unestablished.
These findings concern different measurement boundaries and cannot be pooled
into a new favorable estimator.

The contribution is a derivation and implementation analysis with explicit
negative computational evidence. We provide the side-specific proximal formula,
the contraction-based box reduction, the directional threshold and conjugate,
and the fixed-metric splitting conditions. We compare exact objective-matched
references and retain compatibility limits, unsuccessful observations and
original-space acceptance. Component comparisons examine formulation choice,
conditioning and workspace reuse while stating where no causal intervention
was performed. The results establish mathematical equivalence and tested
implementation behavior within their stated domains, not a uniform runtime
advantage, operational safety or a new expressibility theorem beyond the
identified prior art.

The following sections position the methods within convex proximal and
quadratic-programming literature, define the model, derive each proposed
update separately and state the theory's limits. Experimental design and
results keep the two populations and confidence procedures distinct. The
ablation analysis is explicitly limited to retained component comparisons;
discussion and limitations explain why correct original-row updates need
not improve complete-operation time. Reproducibility relies on the exact
archived inputs and outputs rather than on post-hoc reconstruction of a
successful subset.
